\chapter{Conclusão}

\section{Introdução}

Este trabalho explorou a aplicação do Modelo de Atores em um contexto não tradicional: aplicações centralizadas desenvolvidas em Rust. Através da implementação de uma versão baseada em atores do Patch-Hub, foi possível avaliar os benefícios e desafios de aplicar este padrão arquitetural fora de seu domínio original de sistemas distribuídos.

\section{Principais Contribuições}

Este trabalho contribuiu para o conhecimento sobre arquitetura de software de várias formas:

\begin{itemize}
    \item \textbf{Proposta Arquitetural}: Desenvolveu uma arquitetura baseada no Modelo de Atores especificamente adaptada para aplicações centralizadas em Rust, incluindo padrões de implementação, convenções de nomenclatura e estratégias de organização de código.
    
    \item \textbf{Implementação Prática}: Demonstrou a viabilidade da abordagem através da implementação de uma aplicação real (Patch-Hub) usando a arquitetura proposta.
    
    \item \textbf{Análise Comparativa}: Forneceu uma análise detalhada das diferenças entre a implementação original monolítica e a versão baseada em atores.
\end{itemize}

\section{Forças da Abordagem}

% PLACEHOLDER: This section is marked as incomplete in Index.md
% TODO: Add content about strengths of the actor model approach
% This section should cover:
% - Key benefits observed during implementation
% - Advantages over the original monolithic approach
% - Specific improvements in areas like testability, maintainability, etc.

\subsection{Isolamento e Concorrência}

% PLACEHOLDER: Content to be developed
% Discuss how the actor model improved isolation between components
% Explain the benefits of natural concurrency through message passing
% Analyze the impact on system responsiveness and user experience

\subsection{Testabilidade}

% PLACEHOLDER: Content to be developed
% Describe how the actor model improved testing capabilities
% Discuss the benefits of mock implementations and dependency injection
% Explain how testing became more focused and reliable

\subsection{Manutenibilidade}

% PLACEHOLDER: Content to be developed
% Analyze how the actor model improved code organization
% Discuss the benefits of clear separation of concerns
% Explain how the architecture made the system easier to understand and modify

\section{Fraquezas e Limitações}

% PLACEHOLDER: This section is marked as incomplete in Index.md
% TODO: Add content about weaknesses and limitations
% This section should cover:
% - Challenges encountered during implementation
% - Performance overhead of the actor model
% - Complexity introduced by message passing
% - Limitations of applying distributed systems patterns to centralized applications

\subsection{Overhead de Performance}

% PLACEHOLDER: Content to be developed
% Discuss the performance cost of message passing
% Analyze memory usage implications
% Compare execution overhead with the original implementation

\subsection{Complexidade de Desenvolvimento}

% PLACEHOLDER: Content to be developed
% Discuss the learning curve for developers
% Explain the complexity introduced by message-based communication
% Analyze the impact on development time and debugging

\subsection{Limitações do Contexto}

% PLACEHOLDER: Content to be developed
% Discuss limitations of applying distributed patterns to centralized applications
% Explain scenarios where the actor model might not be suitable
% Analyze trade-offs specific to the centralized application domain

\section{Vale a Pena?}

% PLACEHOLDER: This section is marked as incomplete in Index.md
% TODO: Add content evaluating whether the approach is worth it
% This section should cover:
% - Overall assessment of benefits vs. costs
% - Recommendations for when to use this approach
% - Criteria for evaluating architectural pattern suitability
% - Lessons learned and practical implications

\subsection{Avaliação Geral}

% PLACEHOLDER: Content to be developed
% Provide an overall assessment of the actor model approach
% Weigh the benefits against the costs and complexity
% Make recommendations based on the findings

\subsection{Critérios de Aplicabilidade}

% PLACEHOLDER: Content to be developed
% Define criteria for when to apply the actor model to centralized applications
% Identify characteristics of applications that would benefit from this approach
% Provide guidance for architectural decision-making

\subsection{Lições Aprendidas}

% PLACEHOLDER: Content to be developed
% Summarize key lessons learned during the implementation
% Provide insights for future similar projects
% Discuss practical considerations for developers

\section{Trabalhos Futuros}

% PLACEHOLDER: This section is marked as incomplete in Index.md
% TODO: Add content about future works
% This section should cover:
% - Potential extensions to the current work
% - Areas for further research
% - Improvements to the proposed architecture
% - Applications to other domains or languages

\subsection{Melhorias na Arquitetura}

% PLACEHOLDER: Content to be developed
% Discuss potential improvements to the proposed architecture
% Identify areas where the current implementation could be enhanced
% Propose new patterns or conventions based on lessons learned

\subsection{Extensões de Funcionalidades}

% PLACEHOLDER: Content to be developed
% Discuss additional features that could be implemented
% Explore how the actor model could support more complex scenarios
% Identify opportunities for expanding the current implementation

\subsection{Aplicações em Outros Domínios}

% PLACEHOLDER: Content to be developed
% Explore other domains where this approach might be beneficial
% Discuss potential applications in different types of applications
% Identify research opportunities in related areas

\subsection{Ferramentas e Frameworks}

% PLACEHOLDER: Content to be developed
% Discuss the potential for creating reusable frameworks
% Explore tooling that could support actor-based development in Rust
% Identify opportunities for ecosystem development

\section{Considerações Finais}

Este trabalho demonstrou que é possível aplicar o Modelo de Atores com sucesso em aplicações centralizadas, embora com trade-offs significativos. A arquitetura proposta oferece benefícios claros em termos de testabilidade, manutenibilidade e extensibilidade, mas introduz complexidade adicional que deve ser cuidadosamente considerada.

A escolha de aplicar padrões arquiteturais deve sempre ser baseada em uma análise cuidadosa dos requisitos específicos do projeto, dos trade-offs envolvidos e do contexto de desenvolvimento. Embora o Modelo de Atores possa não ser adequado para todos os tipos de aplicações centralizadas, ele representa uma ferramenta valiosa no arsenal arquitetural quando aplicado apropriadamente.

O conhecimento gerado por este trabalho contribui para a compreensão mais ampla de como padrões arquiteturais podem ser adaptados e aplicados em contextos não tradicionais, abrindo caminho para futuras explorações e inovações na área de arquitetura de software.

